% Imporditud: tools/import_eksperimendid.py
% Allikas: Eesti füüsikaolümpiaadi eksperimendi ülesannete
%   kogu (Rael Kalda, 2023), ülesanne Ü84, lahendus L84.
\setAuthor{EFO žürii}
\setRound{lõppvoor}
\setYear{2001}
\setNumber{G E1}
\setDifficulty{3}
\setTopic{elektriahelad}
\setVahendid{taskulambipirn, taskulambipatarei, reostaat, lüliti, voltmeeter, ampermeeter (kasutada ampermeetrina testrit mõõtepiirkonnas 200 mA), juhtmed, millimeeterpaber}
\setPunktid{10}
\setAllikas{Eksperimendi kogu Ü84, lahendus lk 69}
\setLahendusseis{toores}

\prob{Taskulambipirni takistus 2}
Määrake taskulambipirni takistus toatemperatuuril (testri oommeetrina kasutamise eest punkte ei anta).

\hint

\solu
Vihje: Et taskulambipatarei pinget ei saa muuta, siis kasutage reguleeritava pinge saamiseks reostaati potentsiomeetrilises lülituses.

Lahendus: Hõõglambi takistus sõltub hõõgniidi temperatuurist. Temperatuuri tõustes takistus suureneb. Hõõgniidi temperatuur sõltub voolutugevusest. Mida suurem on voolutugevus, seda suurem on temperatuur. Lambi hõõgniit on toatemperatuuril siis, kui voolutugevus lambis on 0 ehk pinge lambi hõõgniidi otstel võrdub nulliga. Mõõdame voolutugevuse lambis pinge erinevate väärtuste korral ja arvutame igal pinge väärtusele vastava takistuse väärtuse. Erilist tähelepanu tuleb pöörata mõõtmistele pinge väikestel väärtustel. Joonistame graafiku, mille ühel teljel on pinge, teisel --- takistuse väärtus. Pikendades graafiku väärtuseni U = 0, saame hõõgniidi takistuse toatemperatuuril. Pinge muutmiseks kasutame reostaati pingejagajana. Vooluringi skeem on toodud järgmisel joonisel:

A V
\probend
